% Physical unit 006; accepted Session 02.
\GordanUnitStart{006}{printed page 1}
\begin{center}
{\large CAPUT I.}\\[0.8em]
\S\ 1.
\end{center}

Sphaerois est superficies, rotatione ellipseos circa minorem axem \gq{$b$} formata. Quo in motu singula ellipseos peripheriae puncta describunt circulos (parallelos), quorum maximus ab extremis punctis axis majoris \gq{$a$} descriptus, aequator appellatur. ---

\begin{center}\S\ 2.\end{center}
Axis minor ellipseos, circa quem facta est rotatio, nominatur axis sphaeroidis, ejus extrema puncta, poli, omne autem planum, illum axem continens, meridianus.

\begin{center}\S\ 3.\end{center}
Angulus \gq{$\psi$} is, qui efficitur meridiano, continenti punctum \gq{$A$}, et (quovis) primo meridiano, appellatur longitudo; angulus \gq{$B$} autem, quem normalis in $A$ in sphaeroide erecta cum plano aequatoris effingit, latitudo puncti $A$. Inde vides, omne punctum sphaeroidis longitudine et latitudine determinari. ---

\begin{center}\S\ 4.\end{center}
Omnia puncta ejusdem longitudinis $\psi$, cum in eodem meridiano, et omnia puncta ejusdem latitudinis $B$ in eodem circulo sita sint, illum angulum $\psi$ nominamus longitudinem meridiani, eadem puncta continentis, hunc angulum $B$ latitudinem circuli, eadem puncta continentis. ---

\begin{center}\S\ 5.\end{center}
Praeter latitudinem $B$ utamur ad situm circuli determinandum alio angulo \gq{$l$}. Nam, cum, radii loco sphaeroidis, axi \gq{$a$} utentes, circa centrum sphaeram describimus, planum circuli sphaeroidis hanc quoque circulo secat. Linearum normalium, in punctis periphe-
\GordanBottomFolio{1}
